% Настройки абзацного отступа и списков
\setlength{\parindent}{1.25cm}

% Глобальная настройка всех списков (enumitem)
\setlist{
	leftmargin=0pt,
	nosep,
	wide=\parindent,
	labelsep=0.5em,
	listparindent=\parindent
}
\setlist[itemize,1]{label=--} 
\setlist[enumerate,1]{label=\arabic*.}

% Настройка заголовков через KOMA-Script (вместо titlesec)
\RedeclareSectionCommand[beforeskip=12pt, afterskip=6pt, indent=1.25cm]{section}
\RedeclareSectionCommand[beforeskip=12pt, afterskip=6pt, indent=1.25cm]{subsection}

% Команда против переносов
\newcommand{\nohyphens}{\hyphenpenalty=10000\exhyphenpenalty=10000\sloppy}

% Настройка подписей (caption)
\captionsetup{
	justification=centering,
	singlelinecheck=false
}
\DeclareCaptionLabelFormat{simple}{#1~#2}
\captionsetup{labelformat=simple}

% Настройка ссылок
\urlstyle{same}

% Настройка таблиц
\DeclareCaptionLabelSeparator{emdash}{ --- }
\captionsetup[table]{
	labelsep=emdash,
	justification=raggedright,
	singlelinecheck=false,
	position=top,
	skip=6pt
}

% Стили для листингов (JSON и TypeScript)
\lstdefinelanguage{json}{
	basicstyle=\ttfamily\footnotesize,
	numbers=none,
	breaklines=true,
	stringstyle=\color{blue},
}

\lstdefinelanguage{TypeScript}{
	keywords={export, async, function, return, await, string, const, let},
	keywordstyle=\color{blue}\bfseries,
	comment=[l]{//},
	morecomment=[s]{/*}{*/},
	stringstyle=\color{red},
	morestring=[b]',
	morestring=[b]",
	morestring=[b]{`}
}

% Настройка библиографии под СФУ
\DefineBibliographyStrings{russian}{
	urlseen = {дата обращения:},
}